\exercisegroup

\begin{enumerate}
\def\labelenumi{\arabic{enumi})}
\item
  设 \(E_{0}=Q_{p}^{N}\) 是 p 进域 \(Q_{p}\)（GT, III, § 6, 习题 23）上的向量空间，它是由每个都恒等于 \(Q_{p}\) 的因子所组成的可数无穷乘积。设 \(P\subset E_{0}\) 是集合 \(Z_{p}^{N}\)，并设 E 是由 P 生成的 \(E_{0}\) 的向量子空间。在加群 P 上我们考虑诸因子 \(Z_{p}\) 的拓扑的乘积紧拓扑，并用 B 表示 P 关于此拓扑的 0 的邻域滤子。证明：B 是 E 上某个与 E 的加群结构相容的拓扑 T 的 0 的基本邻域系，该拓扑满足 \((\mathrm{EVT}_{\mathrm{I}}^{\prime})\) 与 \((\mathrm{EVT}_{\mathrm{III}}^{\prime})\) 但不满足 \((\mathrm{EVT}_{\mathrm{II}}^{\prime})\)（证明位似 \(x\mapsto x/p\) 在 E 中不连续）。
\item
  设 K 是非离散拓扑除环，\(K_{0}\) 是赋以离散拓扑的除环 K。\(K_{0}\) 上的离散拓扑与其加群结构相容，且当我们把 \(K_{0}\) 看作 K 上的向量空间时，它满足公理 (EVT \(_{II}\)) 与 (EVT \(_{III}\)) 但不满足 (EVT \(_{I}\))。
\item
  对每个实数 \(\alpha > 0\)，设 \(G_{\alpha}\) 是拓扑群 \(\mathbf{R} / \alpha \mathbf{Z}\)，并设 \(G\) 是拓扑乘积群 \(\prod_{\alpha} G_{\alpha}\)（\(\alpha\) 取遍数 \(>0\) 组成的集合）。对每个 \(x \in \mathbf{R}\)，
\end{enumerate}

设 \(t_{\alpha}(x)\) 是 \(x\) 在 \(G_{\alpha}\) 中的典范像；映射 \(\phi : x \to (t_{\alpha}(x))\) 是 \(\mathbf{R}\) 到 \(G\) 中的一个连续单射同态。我们在 \(\mathbf{R}\) 上考虑由 \(\phi\) 取的 \(G\) 的拓扑的原像拓扑，并用 \(E\) 表示由赋以此拓扑的 \(\mathbf{R}\) 所成的拓扑群。证明：当把 \(E\) 看作 \(\mathbf{R}\) 上的向量空间时，其拓扑满足 \((\mathrm{EVT}_{I}^{\prime})\) 与 \((\mathrm{EVT}_{II}^{\prime})\) 但不满足 \((\mathrm{EVT}_{III}^{\prime})\)。

\begin{enumerate}
\def\labelenumi{\arabic{enumi})}
\setcounter{enumi}{3}
\item
  设 E 是带赋值的除环 K 上的一个向量空间；假设 E 带有与其加群结构相容的一个可度量化拓扑。进一步假设此拓扑满足公理 (EVT \(_{I}\)) 与 (EVT \(_{II}\))；证明：若两个可度量化群 K, E 中有一个是完备的，则该拓扑也满足 (EVT \(_{III}\))，从而与 E 的向量空间结构相容（cf.~GT, IX, § 5, 习题 23）。
\item
  设 \(\mathbf{K}\) 是非离散赋值域，\(S\) 是任意一个无穷集合。
\end{enumerate}

\begin{enumerate}
\def\labelenumi{\alph{enumi})}
\item
  设 \(\mathbf{D} = (a_n)\) 是 \(\mathbf{S}\) 的元素的一个可数无穷族。对每个 \(\lambda \in \mathbf{K}\)（满足 \(|\lambda| \leqslant 1\)），设 \(u_{\lambda}\) 是赋范空间 \(\mathcal{B}_{\mathbf{K}}(\mathbf{S})\)（I, 第 4 页, 例；该空间由 \(\mathbf{S}\) 到 \(\mathbf{K}\) 中的有界映射组成）的元素，满足 \(u_{\lambda}(a_n) = \lambda^n\)（对每个 \(n \in \mathbf{N}\)），且 \(u_{\lambda}(b) = 0\)（对每个 \(b \notin \mathbf{D}\)）。证明族 \((u_{\lambda})\) 是代数无关的。
\item
  推出向量空间 \(\mathcal{B}_{\mathbf{K}}(\mathbf{S})\) 的每个基与 \(\mathrm{K}^{\mathrm{s}}\) 有相同的基数（利用 \(a\)），证明 \(\mathcal{B}_{\mathbf{K}}(\mathbf{S})\) 的每个基的基数至少等于 Card (K)；另一方面注意 \(\operatorname{Card}(\mathcal{B}_{\mathbf{K}}(\mathbf{S})) = \operatorname{Card}(\mathbf{K}^{\mathrm{s}})\) 并利用 A, II, § 2, 习题 22）。
\item
  以同样的方式证明向量空间 \(\ell_{\mathbf{K}}^{1}(\mathbf{S})\) 的每个基具有 \((\mathbf{K}\times \mathbf{S})^{\mathbf{N}}\) 的基数。
\end{enumerate}

\begin{enumerate}
\def\labelenumi{\arabic{enumi})}
\setcounter{enumi}{5}
\tightlist
\item
  设 \(\mathbf{K}\) 是非离散赋值除环。证明：对空间 \(\ell_{\mathbf{K}}^{1}(\mathbf{N})\)（由形如 \(x = (\xi_n)\) 的、由 \(\mathbf{K}\) 中元素组成的绝对可和序列构成）而言，范数 \(\| x \|_1 = \sum_{n=0}^{\infty} |\xi_n|\) 与 \(\| x \| = \sup_n |\xi_n|\) 不等价（cf.~GT, IX, § 3.3, prop. 7）；证明 \(\ell_{\mathbf{K}}^{1}(\mathbf{N})\)（带范数 \(\| x \|\)）从不完备，即使 \(\mathbf{K}\) 是完备的；它在 \(\mathcal{B}_{\mathbf{K}}(\mathbf{N})\) 中的闭包是什么？
\end{enumerate}

¶ 7) * 设 A 是一个离散赋值环，v 是 A 的分式除环 K 的赋范赋值；在 K 上取绝对值 \(a^{v}\)，其中 0 \textless{} a \textless{} 1。设 E 是 K 上的一个赋范向量空间，其范数满足超度量不等式

\[
\| x + y \| \leqslant \sup (\| x \|, \| y \|).
\]

\begin{enumerate}
\def\labelenumi{\alph{enumi})}
\item
  用 M 表示由那些 \(x \in E\)（满足 \(\| x \| \leqslant 1\)）所组成的集合，用 \(\pi\) 表示 A 的一个单值化元；M 是一个 A-模，而 \(M / \pi M\) 是 A 的剩余除环 \(k = A / \pi A\) 上的一个向量空间。设 \((e_{\lambda})_{\lambda \in L}\) 是 M 的元素的一个族，使得诸 \(e_{\lambda}\) 在 \(M / \pi M\) 中的像构成这个 \(k\) -向量空间的一个基。证明 \((e_{\lambda})\) 是 E 中的一个无关族，且由 \((e_{\lambda})\) 生成的 E 的向量子空间 F 在 E 中稠密。
\item
  若置 \(\| x\| _1 = \sup_{\lambda}|\xi_\lambda |,\)（对 F 中每个 \(x = \sum_{\lambda}\xi_{\lambda}e_{\lambda}\)），证明在 F 上范数 \(\| x\|\) 与 \(\| x\| _1\) 等价。
\item
  设 \(\mathbf{K}\) 是完备的。由 \(a)\) 与 \(b)\) 推出：若 \(\mathbf{L}\) 有限，则完备化 \(\hat{\mathbf{E}}\)（即 \(\mathbf{E}\) 的完备化）同构于 \(\mathbf{K}^{\mathrm{L}}\)；若 \(\mathbf{L}\) 无限，则 \(\hat{\mathbf{E}}\) 同构于子空间 \(\mathcal{C}_{\mathbf{K}}^{0}(\mathbf{L})\)，即 \(\mathcal{B}_{\mathbf{K}}(\mathbf{L})\) 中由那些族 \((\xi_{\lambda})\)（它们满足 \(\lim \xi_{\lambda} = 0\)，关于 \(\mathbf{L}\) 的有限子集的补集所组成的滤子）所组成的子空间。
\item
  假设 K 与 E 完备；设 G 是 K 上另一个完备赋范空间，其范数满足超度量不等式。证明：在（如有必要）把 \(\mathcal{L}(\mathrm{E};\mathrm{G})\)（GT, X, § 3.2）的范数换成一个等价的范数之后，\(\mathcal{L}(\mathrm{E};\mathrm{G})\) 等距同构于由 G 中元素组成的诸族 \((y_{\lambda})_{\lambda \in L}\)（它们满足 \(\sup_{\lambda \in L}\| y_{\lambda}\| < +\infty\)）所成的向量空间，后者带范数 \(\sup_{\lambda \in L}\| y_{\lambda}\|\)（它也是一个超度量范数）。
\end{enumerate}

\begin{enumerate}
\def\labelenumi{\arabic{enumi})}
\setcounter{enumi}{7}
\item
  设 E 是非离散拓扑除环 K 上的一个拓扑向量空间。为使存在点 \((0,0)\)（在 \(K \times E\) 中）的一个邻域，使映射 \((\lambda, x) \mapsto \lambda x\) 在此邻域上一致连续，必须且只需存在 E 中 0 的一个邻域 \(V_{0}\)，使得诸集合 \(\lambda V_{0}\)（其中 \(\lambda\) 取遍 K 中 \(\neq 0\) 的元素）构成 E 中 0 的一个基本邻域系。当 K 是非离散赋值除环且 E 是 Hausdorff 的时，证明 E 的一致结构此时是可度量化的。
\item
  把 I, 第 8 页 的 prop. 5 推广到如下的情形：诸空间 \(\mathrm{E}_i\)（\(\leqslant i \leqslant n\)）与 \(\mathrm{F}\) 是任意非离散拓扑域上的拓扑向量空间。
\item
  设 E 是非离散赋值除环 K 上的一个完备 Hausdorff 拓扑向量空间。用 F 表示 E 的一个向量子空间，用 T 表示 F 上由 E 的拓扑 \(T'\) 诱导的拓扑；设 B 是拓扑 T 的 0 的闭均衡邻域的一个基本系。设 \(F_{0}\) 是 E 的向量子空间，由闭包 \(\overline{V}\)（在 E 中关于 \(T'\) 取的、诸集合 \(V \in B\) 的闭包）生成；诸集合 \(\overline{V}\) 构成某个拓扑 \(T_{0}\)（\(F_{0}\) 上的拓扑）的 0 的一个基本邻域系，该拓扑与 \(F_{0}\) 的向量空间结构相容；对此拓扑，\(F_{0}\) 是完备的，且 \(T_{0}\) 在 F 上诱导的拓扑与 T 恒同。
\item
  在非离散拓扑除环 K 上的拓扑向量空间 E 中，存在 0 的闭邻域的一个基本系 B，满足条件 (EV \(_{II}\)) 与 (EV \(_{III}\)) 以及下面两个：
\end{enumerate}

\((\mathrm{EV}_{\mathrm{Ia}})\) 对每个 \(\mathbf{V} \in \mathfrak{B}\)，存在 \(\mathbf{W} \in \mathfrak{B}\) 以及 0 的一个邻域 \(\mathbf{U}\)（\(\mathbf{K}\) 中的邻域），使得 \(\mathrm{UW} \subset \mathrm{V}\)。\((\mathrm{EV}_{\mathrm{Ib}})\) 对每个 \(x \in \mathbf{E}\) 及每个 \(\mathbf{V} \in \mathfrak{B}\)，存在 \(\lambda \neq 0\)（\(\mathbf{K}\) 中的元素），使得 \(\lambda x \in \mathbf{V}\)。反之，设 \(\mathbf{E}\) 是 \(\mathbf{K}\) 上的一个向量空间，并设 \(\mathfrak{B}\) 是 \(\mathbf{E}\) 上满足条件 \((\mathrm{EV}_{\mathrm{Ia}}), (\mathrm{EV}_{\mathrm{Ib}}), (\mathrm{EV}_{\mathrm{II}})\) 与 \((\mathrm{EV}_{\mathrm{III}})\) 的一个滤子基。证明 \(\mathbf{E}\) 上存在（且仅存在一个）与 \(\mathbf{E}\) 的向量空间结构相容的拓扑，使 \(\mathfrak{B}\) 是它的 0 的基本邻域系。

\begin{enumerate}
\def\labelenumi{\arabic{enumi})}
\setcounter{enumi}{11}
\item
  设 K 是离散域，E 是 K 上两个未定元的形式级数环 A = K{[}{[}X, Y{]}{]}（A, IV, 第 36 页）的分式除环。对每个 \(n \geqslant 0\)，设 \(V_{n} \subset A\) 是总次数至少等于 n 的形式级数所成的集合。证明：在 E 中，诸集合 \(V_{n}\) 构成 0 的一个基本邻域系，对应于与 E（在 K 上）的向量空间结构相容的一个拓扑，对此拓扑 E 是可度量化且完备的；若进一步 K 是有限域，则 E 是局部紧的。证明 E × E 到 E 中的 K-双线性映射 \((u, v) \mapsto uv\) 在点 \((0, 0)\) 处连续，但存在 \(u_{0} \in E\) 使 \(v \mapsto u_{0}v\) 在 E 中不连续（例如 \(u_{0} = 1/X\)）。
\item
  设 E 是 R 上无穷维的一个向量空间，并设 T 是 E 的所有吸收集与均衡集所成的族。证明 T 不满足公理 \((\mathrm{EV}_{\mathrm{III}})\)（换言之，不是与 E 的加群结构相容的某个拓扑的 0 的基本邻域系）。为此，考虑 E 中一个无穷无关族 \((e_{n})_{n\geqslant1}\)；对每个整数 \(n\geqslant1\)，设 \(A_{n}\) 是形如 \(\sum_{i=1}^{n}t_{i}e_{i}\) 的点（满足 \(|t_{i}|\leqslant1/n\)，对 \(1\leqslant i\leqslant n\)）所成的集合；设 A 是诸 \(A_{n}\) 的并，V 是与 E 中由诸 \(e_{n}\) 生成的子空间互补的一个子空间，并用 C 表示集合 \(A+V\)；证明不存在集合 \(M\in T\) 使 \(M+M\subset C\)。
\end{enumerate}

¶ 14) 设 K 是 Hausdorff 拓扑除环，\((E_{t})_{t\in I}\) 是 K 上的 Hausdorff 拓扑向量空间的一个无穷族，其中没有一个空间是单点 0。我们在 \(F = \prod_{t\in I} E_t\) 上考虑拓扑 \(\mathcal{T}\)，它与 \(F\) 的加群结构相容，其 0 的一个基本邻域系由诸乘积 \(\prod_{\iota \in I} V_{\iota}\) 组成，其中对每个 \(\iota \in I\)，集合 \(V_{t}\) 是 E 中 0 的一个邻域（这一拓扑严格细于乘积拓扑；cf.~GT, III, § 2, 习题 23）。用 \(T_{0}\) 表示 T 在 F 的子空间 \(E = \bigoplus_{i \in I} E_{i}\) 上诱导的拓扑；对于拓扑 T，E 在 F 中是闭的，且若诸 \(E_{t}\) 中每个都完备，则 F 关于拓扑 T 是完备的，因此 E 关于拓扑 \(T_{0}\) 是完备的（GT, III, § 3, 习题 10）。

\begin{enumerate}
\def\labelenumi{\alph{enumi})}
\item
  证明：若 K 中存在 0 的一个右有界邻域（GT, III, § 6, 习题 12）（特别地，若 K 是赋值除环），则拓扑 \(\mathcal{T}_0\) 与 E 的向量空间结构相容。若进一步 K 不是离散的，则对于任何细于 \(\mathcal{T}_0\) 且与 E 的向量空间结构相容的拓扑，E 都不是 Baire 空间。b) 此外，若 K 中不存在 0 的任何右有界邻域（参见 c)），给出一个族 \((\mathrm{E}_t)\) 的例子，使拓扑 \(\mathcal{T}_0\) 不与 E 的向量空间结构相容。
\item
  设 \(\mathbf{A} = \mathbf{R}[\mathbf{X}]\) 是 \(\mathbf{R}\) 上一个未定元的多项式环。对每个序列 \(s = (\varepsilon_n)_{n\geqslant 0}\)（由实数 \(>0\) 组成的序列），用 \(\mathrm{V}_s\) 表示多项式 \(\sum_k a_k\mathrm{X}^k\in \mathrm{A}\)（满足 \(|a_k| < \varepsilon_k\)，对每个 k）所成的集合。设 T 是诸 \(V_{s}\) 所成的集合，其中 s 取遍由数 \textgreater0 组成的序列所成的集合。证明 T 是与 A 的环结构相容的某个拓扑的 0 的对称邻域的一个基本系。设 \(\mathbf{K} = \mathbf{R}(\mathbf{X})\) 是 A 的分式除环；用 S 表示 K 的形如 \(\mathrm{U}(1 + \mathrm{U})^{-1}\) 的子集所成的族，其中 U 取遍不含 1 的诸 \(V_{s}\)；证明 S 是与 K 的除环结构相容的某个拓扑的 0 的一个基本邻域系，且 K 中不存在有界的 0 邻域。
\end{enumerate}

\(d\) ) 对每个 Hausdorff 拓扑除环 \(\mathbf{K}\)，证明存在一个集合 I，使得在 \(\mathbf{F} = \mathbf{K}^{\mathrm{I}}\) 上，前面定义的拓扑 \(\mathcal{T}\) 不与 F 的向量空间结构相容。
% semantic-fragments: legacy-input-seam
\relax{}
